\documentclass[11pt,a4paper]{article}

\usepackage[margin=2.5cm]{geometry}
\usepackage{listings}
\usepackage{xcolor}
\usepackage{hyperref}
\usepackage{titlesec}
\usepackage{parskip}
\usepackage{booktabs}
\usepackage{microtype}
\usepackage{fancyhdr}
\usepackage{caption}

% ── Colour palette ────────────────────────────────────────────────────────────
\definecolor{codebg}{HTML}{F5F5F5}
\definecolor{codeframe}{HTML}{CCCCCC}
\definecolor{kwcolor}{HTML}{0000CC}
\definecolor{strcolor}{HTML}{008800}
\definecolor{cmtcolor}{HTML}{888888}
\definecolor{numcolor}{HTML}{880000}

% ── Go listing style ──────────────────────────────────────────────────────────
\lstdefinelanguage{Go}{
  morekeywords={break,case,chan,const,continue,default,defer,else,fallthrough,
    for,func,go,goto,if,import,interface,map,package,range,return,select,
    struct,switch,type,var,nil,true,false,make,new,len,cap,append,close,
    delete,copy,error,string,int,int64,float64,bool,byte},
  sensitive=true,
  morecomment=[l]{//},
  morecomment=[s]{/*}{*/},
  morestring=[b]",
  morestring=[b]`,
}

\lstset{
  language=Go,
  backgroundcolor=\color{codebg},
  frame=single,
  rulecolor=\color{codeframe},
  basicstyle=\ttfamily\footnotesize,
  keywordstyle=\color{kwcolor}\bfseries,
  stringstyle=\color{strcolor},
  commentstyle=\color{cmtcolor}\itshape,
  numberstyle=\tiny\color{numcolor},
  numbers=left,
  numbersep=8pt,
  breaklines=true,
  breakatwhitespace=false,
  tabsize=4,
  showstringspaces=false,
  captionpos=b,
}

% ── Header / footer ───────────────────────────────────────────────────────────
\pagestyle{fancy}
\fancyhf{}
\rhead{CSU33032 Advanced Networks}
\lhead{HTTP Web Proxy}
\cfoot{\thepage}

% ── Section spacing ───────────────────────────────────────────────────────────
\titlespacing*{\section}{0pt}{1.5ex plus .5ex}{0.8ex}
\titlespacing*{\subsection}{0pt}{1ex plus .3ex}{0.5ex}

% ─────────────────────────────────────────────────────────────────────────────
\begin{document}

\begin{titlepage}
  \centering
  \vspace*{3cm}
  {\Large\bfseries CSU33032 Advanced Networks\par}
  \vspace{0.5cm}
  {\huge\bfseries HTTP Web Proxy\par}
  \vspace{0.5cm}
  {\large Protocol Design and Implementation Report\par}
  \vspace{2cm}
  {\large Liangzicheng\par}
  \vspace{0.5cm}
  {\large\today\par}
  \vfill
  {\small Trinity College Dublin\par}
\end{titlepage}

% ─────────────────────────────────────────────────────────────────────────────
\tableofcontents
\newpage

% =============================================================================
\section{Overview}
% =============================================================================

This report describes the design and implementation of an HTTP/HTTPS web proxy
written in Go. The proxy listens on TCP port 4000, accepts browser connections,
and forwards requests to origin servers. Three features are implemented beyond
basic forwarding: response caching, host-based site blocking, and a live
terminal user interface (TUI) for monitoring and management.

% =============================================================================
\section{Protocol Design}
% =============================================================================

\subsection{Supported HTTP Methods}

The proxy handles two classes of traffic:

\begin{itemize}
  \item \textbf{Plain HTTP} (\texttt{GET}, \texttt{HEAD}, and others): the proxy
    receives the full request, optionally serves it from cache, or forwards it
    to the origin server and returns the response.
  \item \textbf{HTTPS tunnelling via \texttt{CONNECT}}: the proxy establishes a
    raw TCP tunnel to the destination host and splices the client and server
    streams together. Because the traffic is encrypted end-to-end, it is not
    cached.
\end{itemize}

\subsection{Caching}

Responses to \texttt{GET} and \texttt{HEAD} requests are cached in memory.
The cache key is the concatenation \texttt{method|host|path}, giving
per-resource granularity. On a cache hit the stored byte slice is written
directly to the client socket without contacting the origin server.

The current implementation is a simple in-memory store with no expiry (no TTL
or \texttt{Cache-Control} header parsing). This is intentional for the scope
of this assignment; a production proxy would honour RFC 7234.

\subsection{Site Blocking}

Before any network I/O the proxy consults a \texttt{map[string]bool} of blocked
hostnames. If the destination host (port stripped) is present, the proxy returns
\texttt{HTTP/1.1 403 Forbidden} immediately. The map is checked under a mutex so
it can be updated at runtime from the management console without data races.

\subsection{Concurrency Model}

Each accepted TCP connection is handed off to a new goroutine. Shared state
(\texttt{blockedSites}, \texttt{cache}, \texttt{logs}) lives in a single
\texttt{proxyData} struct protected by a \texttt{sync.Mutex}. The lock is held
only for map reads and writes, keeping the critical section short.

% =============================================================================
\section{System Architecture}
% =============================================================================

The codebase comprises two source files:

\begin{itemize}
  \item \textbf{\texttt{proxy.go}} --- TCP listener, connection handler, cache,
    and blocking logic.
  \item \textbf{\texttt{main.go}} --- Terminal UI (Bubbletea framework),
    statistics aggregation, and management console.
\end{itemize}

Communication between the proxy goroutines and the TUI follows the
\emph{Elm Architecture}: goroutines emit typed messages (\texttt{RequestMsg},
\texttt{LogMsg}) via \texttt{program.Send()}, and the TUI's \texttt{Update()}
function processes them on its own goroutine. This decouples latency-sensitive
I/O from rendering and avoids holding any lock during UI updates.

\begin{figure}[h]
\centering
\begin{verbatim}
  Browser
     |
  TCP :4000
     |
  listener (startProxy)
     |
  goroutine per connection (handleConnection)
     |-- blocked? --> 403 Forbidden
     |-- CONNECT  --> TCP tunnel (io.Copy both directions)
     |-- GET/HEAD cached? --> write from cache
     |-- GET/HEAD fresh  --> dial origin, cache response
     |
  program.Send(RequestMsg / LogMsg)
     |
  Bubbletea Update() --> View() --> terminal
\end{verbatim}
\caption{Request flow through the proxy.}
\end{figure}

% =============================================================================
\section{Implementation}
% =============================================================================

\subsection{Proxy Core (\texttt{proxy.go})}

\subsubsection{Shared State}

\begin{lstlisting}[caption={Shared proxy state structure.}]
type proxyData struct {
    mu           sync.Mutex
    blockedSites map[string]bool   // hostname -> blocked
    cache        map[string][]byte // "method|host|path" -> raw HTTP response
    logs         map[string]string // reserved for future persistent logging
}
\end{lstlisting}

All three maps are initialised in \texttt{main()} and passed by pointer to
both the proxy goroutines and the TUI model, ensuring a single source of truth.

\subsubsection{TCP Listener}

\begin{lstlisting}[caption={Accepting connections and spawning goroutines.}]
func startProxy(sharedData *proxyData, program *tea.Program) {
    listener, err := net.Listen("tcp", ":4000")
    if err != nil {
        log.Fatal("[Proxy Error]", err)
    }
    defer listener.Close()

    for {
        conn, err := listener.Accept()
        if err != nil {
            program.Send(LogMsg{Level: "ERROR", Text: err.Error()})
            continue
        }
        // Each connection runs in its own goroutine; the main loop never blocks.
        go handleConnection(conn, sharedData, program)
    }
}
\end{lstlisting}

\subsubsection{Connection Handler}

The handler performs the following steps in order:

\begin{enumerate}
  \item Parse one HTTP request with \texttt{http.ReadRequest}.
  \item Extract \texttt{Host} and ensure it has a port suffix for dialling.
  \item Check the blocked-sites map under the mutex; return 403 if matched.
  \item For \texttt{CONNECT}: dial the destination, reply 200, and splice
    streams with two \texttt{io.Copy} calls.
  \item For \texttt{GET}/\texttt{HEAD}: check the cache; if hit, write and return.
  \item On a cache miss: dial origin, write the request with
    \texttt{req.Write()}, read the response with \texttt{http.ReadResponse},
    serialise to bytes, forward to client, and store in cache.
\end{enumerate}

\begin{lstlisting}[caption={Blocking check (proxy.go).}]
// Strip port so "example.com:443" matches the key "example.com"
blockHost := host
if h, _, err := net.SplitHostPort(host); err == nil {
    blockHost = h
}
sharedData.mu.Lock()
isBlocked := sharedData.blockedSites[blockHost]
sharedData.mu.Unlock()
if isBlocked {
    clientConn.Write([]byte("HTTP/1.1 403 Forbidden\r\nContent-Length: 0\r\n\r\n"))
    return
}
\end{lstlisting}

\begin{lstlisting}[caption={HTTPS tunnel via CONNECT (proxy.go).}]
if method == "CONNECT" {
    serverConn, err := net.Dial("tcp", dialHost)
    if err != nil {
        clientConn.Write([]byte("HTTP/1.1 502 Bad Gateway\r\nContent-Length: 0\r\n\r\n"))
        return
    }
    // Acknowledge the tunnel; browser begins TLS handshake after this.
    clientConn.Write([]byte("HTTP/1.1 200 Connection Established\r\n\r\n"))
    // Splice: browser <-> proxy <-> server, bidirectionally.
    go io.Copy(serverConn, clientConn)
    io.Copy(clientConn, serverConn)
    serverConn.Close()
    return
}
\end{lstlisting}

\begin{lstlisting}[caption={Cache lookup and hit path (proxy.go).}]
cacheKey := method + "|" + host + "|" + path
if method == "GET" || method == "HEAD" {
    sharedData.mu.Lock()
    cachedResponse, cacheHit := sharedData.cache[cacheKey]
    sharedData.mu.Unlock()

    if cacheHit {
        // Serve the cached raw HTTP response directly — no origin contact.
        clientConn.Write(cachedResponse)
        program.Send(RequestMsg{CacheHit: true, Duration: time.Since(start), ...})
        return
    }
}
\end{lstlisting}

\begin{lstlisting}[caption={Fresh fetch and cache store (proxy.go).}]
serverConn, _ := net.Dial("tcp", dialHost)
defer serverConn.Close()

// req.Write reconstructs a valid HTTP/1.1 request to the origin.
req.Write(serverConn)

// http.ReadResponse respects Content-Length / chunked encoding and returns
// without blocking for connection close (unlike io.ReadAll on the socket).
resp, _ := http.ReadResponse(bufio.NewReader(serverConn), req)
defer resp.Body.Close()

body, _ := io.ReadAll(resp.Body)
resp.Body = io.NopCloser(bytes.NewReader(body))

// Re-serialise full HTTP response (status line + headers + body) into []byte.
var buf bytes.Buffer
resp.Write(&buf)
serverResponse := buf.Bytes()

clientConn.Write(serverResponse)

// Store under mutex; subsequent requests for same resource hit the cache.
sharedData.mu.Lock()
sharedData.cache[cacheKey] = serverResponse
sharedData.mu.Unlock()
\end{lstlisting}

\subsection{Terminal UI (\texttt{main.go})}

The TUI is built with the Bubbletea framework, which implements the Elm
Architecture: a single immutable \texttt{model} is updated by an \texttt{Update}
function in response to messages, and rendered by a \texttt{View} function.

\subsubsection{Message Types}

\begin{lstlisting}[caption={Typed messages sent from proxy goroutines to the TUI.}]
type RequestMsg struct {
    Method   string
    Host     string
    Path     string
    CacheHit bool        // true -> CACHED label
    Blocked  bool        // true -> BLOCKED label
    Duration time.Duration
    Bytes    int
}

type LogMsg struct {
    Level string // INFO | CACHE | BLOCK | ERROR
    Text  string
    Timestamp time.Time
}
\end{lstlisting}

\subsubsection{Four-Tab Interface}

\begin{itemize}
  \item \textbf{Home} --- live statistics: total requests, cache hit rate,
    bytes saved, average latency for cached vs.\ fresh responses.
  \item \textbf{Request Stream} --- scrollable table of every request with
    method, host, path, status (FRESH / CACHED / BLOCKED), and duration.
  \item \textbf{Management Console} --- keyboard-driven site blocking:
    \texttt{i} enters insert mode, \texttt{Enter} adds the site to the shared
    \texttt{blockedSites} map, \texttt{d} removes the last entry.
  \item \textbf{Logs} --- timestamped log entries colour-coded by level.
\end{itemize}

\subsubsection{Runtime Site Blocking via the Console}

\begin{lstlisting}[caption={Adding a blocked site at runtime (main.go).}]
case "enter":
    site := strings.TrimSpace(m.textInput.Value())
    if site != "" {
        m.blockedSites = append(m.blockedSites, site)
        // Write under mutex so proxy goroutines see the change immediately.
        m.proxyShared.mu.Lock()
        m.proxyShared.blockedSites[site] = true
        m.proxyShared.mu.Unlock()
        m.textInput.SetValue("")
    }
\end{lstlisting}

% =============================================================================
\section{Design Choices}
% =============================================================================

\begin{description}

  \item[Goroutine-per-connection.]
    The simplest model for I/O-bound workloads in Go. Each goroutine blocks on
    network I/O independently; the scheduler multiplexes them onto OS threads.
    No explicit thread pool is needed because Go's runtime manages goroutine
    scheduling efficiently.

  \item[\texttt{http.ReadRequest} / \texttt{http.ReadResponse} over raw
    \texttt{io.ReadAll}.]
    These standard-library functions parse framing (Content-Length, chunked
    transfer encoding) and return as soon as the response body is complete.
    Reading raw bytes until EOF would hang on keep-alive connections.

  \item[Shared mutex vs.\ channels.]
    The cache and blocked-sites map are straightforward read-mostly data
    structures. A \texttt{sync.Mutex} with short critical sections is simpler
    and lower-overhead than a dedicated goroutine communicating over channels.

  \item[TUI message passing.]
    Proxy goroutines call \texttt{program.Send()} rather than writing to the
    model directly. This keeps the TUI's state mutations on a single goroutine
    (Bubbletea's event loop), eliminating the need for locking inside the model.

  \item[No TTL / cache eviction.]
    An in-memory cache without expiry is appropriate for a short-lived proxy
    session in an academic context. The trade-off is stale content after server
    updates; a production system would parse \texttt{Cache-Control} and
    \texttt{Expires} headers.

\end{description}

% =============================================================================
\section{Testing and Demonstration}
% =============================================================================

The proxy was tested by configuring Firefox and Chrome to use
\texttt{localhost:4000} as an HTTP proxy. The following was verified:

\begin{itemize}
  \item Plain HTTP sites load correctly via GET forwarding.
  \item HTTPS sites establish tunnels via CONNECT (certificate verification
    occurs end-to-end in the browser, not in the proxy).
  \item Repeated GET requests to the same URL are served from cache with
    significantly lower latency than fresh fetches.
  \item Sites added through the Management Console are immediately blocked;
    the browser receives a 403 response.
  \item The Request Stream and Logs tabs update in real time as traffic flows.
\end{itemize}

% =============================================================================
\section{Conclusion}
% =============================================================================

The proxy correctly forwards HTTP and HTTPS traffic, caches GET/HEAD responses
in memory, and enforces host-based blocking at runtime. The implementation
demonstrates Go's concurrency primitives (goroutines, mutexes) and the
separation of I/O logic from UI state via typed message passing. Future work
would include TTL-based cache expiry, RFC 7234 compliance, persistent block
lists, and HTTPS inspection via a local CA certificate.

\end{document}
